\subsection{向量空间的基}

定理1. 域K上的每个向量空间都是自由K-模。

必须证明每个向量空间都有基；这将由下面更精确的定理得出：

定理2. 给定域K上向量空间E的一个生成系S，以及E中含于S的一个自由子集L，存在E的一个基B，使得 \(L \subset B \subset S\)。

定理1由取 \(\mathbf{L} = \varnothing\) 时的该命题得出。

为证明定理2，我们注意到E中含于S的自由子集的集合 \(\mathfrak{L}\) 按包含关系排序是一个归纳集（集合论，III，§2，第4号），这依据§1第11号；包含L且含于S的自由子集的集合 \(\mathfrak{M}\) 亦是如此。根据Zorn引理，\(\mathfrak{M}\) 有极大元B，且只需证明由B生成的E的向量子空间等于E。这由B的定义及以下引理直接得出：

引理1. 设 \((a_{i})_{i\in I}\) 是 \(\mathbf{E}\) 中元素的一个自由族；若 \(b\in \mathbf{E}\) 不属于子空间 \(\mathbf{F}\)（由 \((a_{i})\) 生成的），则 \(\mathbf{E}\) 的由诸 \(a_{i}\) 与 \(b\) 组成的子集是自由的。

假设有一个关系 \(\mu b + \sum_{i} \lambda_i a_i = 0\)，其中 \(\mu \in K\) 且 \(\lambda_i \in K\)（对所有 \(i \in I\)），族 \((\lambda_i)\) 具有有限支撑；若 \(\mu \neq 0\)，则将得出 \(b = -\sum_{i} (\mu^{-1} \lambda_i) a_i\)，从而 \(b \in F\)，与假设矛盾；因此 \(\mu = 0\)，该关系变为 \(\sum_{i} \lambda_i a_i = 0\)，由假设这蕴含 \(\lambda_i = 0\)（对所有 \(i \in I\)）；由此得引理。

推论. 对向量空间E的子集B，以下性质等价：

\begin{enumerate}
\def\labelenumi{(\alph{enumi})}
\item
  B 是 E 的一个基。
\item
  B 是 E 的一个极大自由子集。
\item
  B 是 E 的一个极小生成系。
\end{enumerate}

这直接由定理 2 得出。

例. 给定一个环 A 和 A 的一个子域 K，A 是 K 上的（右或左）向量空间，因此有一个基；特别地，域 K 的每一个扩域作为 K 上的左（相应地，右）向量空间都有一个基。*因此实数域 R 作为有理数域 Q 上的向量空间有一个（无穷）基；R 的这种基称为哈梅尔基。*

注. 域K上向量空间E的元素族 \((a_{i})_{i\in I}\) 是自由的，当且仅当对所有 \(k\in I\)，\(a_{k}\) 不属于由诸 \(a_{i}\)（下标 \(i\neq k\)）生成的E的任何子空间。我们知道这一条件在任何模中都是必要的（§1，第11号，注1）。由引理1可知它是充分的，这通过反证法并考虑 \((a_{i})\) 的一个极小相关子族立即看出。

